\section{Short square relations and the rank of a pair}
\label{p13:rig:section}

A large intersection of the square-class spaces of two windows
contains many relations, but those relations need not initially have
short support.  We first extract a large subspace with a short basis.
Two independent basis vectors then produce an algebraic curve through
the two window centres.  A resultant argument puts the whole subspace
on one curve, while the relations force many integer points in the
small grid of offsets.  A point count uniform in the degree rules out
this configuration unless the centres lie on a rational alignment.
The distinction between a fixed degree and a growing degree is
essential: the basis supports grow with $L$ throughout the argument.

All estimates in this section are uniform in
$x,y\in I_N$, $|L-\log_2N|\le C$, for fixed $C$, and we put
$B=L+1$.  A logarithm without a subscript is natural.
The conclusion used in the probability argument is
Corollary~\ref{p13:rig:triple}.

\subsection{Square classes on all windows}
\label{p13:rig:notation}

We use the true window
$V_x=\{x-1,x,\ldots,x+L-1\}$, including its left endpoint.
Let $E_x$ be the space of even selections from $V_x$, and let
\[
 \phi_x:E_x\longrightarrow
 \mathbb Q_{>0}^{\times}/\mathbb Q_{>0}^{\times2},
 \qquad I\longmapsto\left[\prod_{n\in I}n\right]
\]
be the square-class map.  Its image is the space $S_x$ already used
above.  Put $d_{xy}=\dim(S_x\cap S_y)$ and, for any finite set of
centres $T$, define the \emph{relative nullity}
\begin{equation}
 \rho_{\rm rel}(T)=\sum_{x\in T}\dim S_x-
                         \dim\left(\sum_{x\in T}S_x\right).
 \label{p13:rig:relative-nullity}
\end{equation}
On the regular set, where every $S_x$ has dimension $L$, this is
the earlier nullity $\rho(T)$.  Formula
\eqref{p13:rig:relative-nullity} also applies to the remaining
windows and separates their local rank defects from dependence
between distinct windows.

A \emph{complete relation} of a pair is a pair of even selections
with the same square class.  We represent its selections by offsets
$I,J\subset\{-1,0,\ldots,L-1\}$ and attach the polynomial
\begin{equation}
 F_w(X,Y)=\prod_{u\in I}(X+u)\prod_{v\in J}(Y+v),
 \qquad w=(I,J).
 \label{p13:rig:word-polynomial}
\end{equation}
Thus $F_w(x,y)$ is an integer square.  The word $w$ is a vector
over $\F_2$; addition means symmetric difference in each block.
The two blocks remain disjoint formal copies even when the actual
windows overlap.  Consequently, independent words give independent
classes of their polynomials in
$\mathbb Q(X,Y)^\times/\mathbb Q(X,Y)^{\times2}$: the factors
$X+u$ and $Y+v$ are distinct irreducible polynomials.

We call $x,y$ \emph{aligned} if there are coprime positive integers
$a,b\le B^3$ such that
\begin{equation}
                   |bx-ay|\le(a+b)B.
 \label{p13:rig:alignment}
\end{equation}
This definition includes nearby centres.  Otherwise they are
nonaligned.  Alignment is a property of the centres; compatibility
of their run events is a separate condition.


\subsection{Extracting short relations}

Primes larger than $10B$ meet each physical block at most once.  Their
parity equations therefore join positions into components; a component
meeting a prime only once is forced to zero.  Local square-product
separation shows that there are only $O_C(B/\log B)$ defective
positions.  Thus the number of free components meeting both blocks is
$d_{xy}+O_C(B/\log B)$; see
Lemmas~\ref{p13:rig:defective} and~\ref{p13:rig:components}.

For medium primes, columns of weight two can be kept as edges of a
graph.  Its short cycles cancel their equations.  Discarding only
columns of weight at least three costs one third of their total
weight, explaining the factor $2/3$ in the following extraction.

\begin{theorem}[A short basis with a controlled loss of dimension]
\label{p13:rig:short-basis}
For every fixed $a\in(0,1)$ there is a subspace $\mathcal K_a$
of complete relations of the pair, with injective projections
onto each physical block and onto the common square class,
such that
\begin{equation}
 \dim\mathcal K_a\ge d_{xy}-\frac23B\log(1/a)
                    -O_{a,C}\!\left(\frac B{\sqrt{\log B}}\right).
 \label{p13:rig:short-dimension}
\end{equation}
It has a basis whose vectors have total physical support at most
\begin{equation}
                    O_{a,C}\!\left(
                         \frac{B^a\log B}{\log\log B}\right).
 \label{p13:rig:short-support}
\end{equation}
Every nonzero vector selects a nonempty even set in each block
and has nonzero common square class.  Neither compatibility nor
nonalignment is assumed.
\end{theorem}

The full extraction proof, including the final small-prime equations
and projection injectivity, is in
Section~\ref{p14:comp:rigidity:short-basis}.

\subsection{One curve for a whole relation space}

Two independent short words first yield a nonzero polynomial through
the window centres.  Proposition~\ref{p13:rig:two-witnesses} constructs
it as the norm of a biquadratic auxiliary function, with degree and
logarithmic height linear in the support size.  Polynomial gcd removal
is performed before specialization, so vanishing at the integer point
does not destroy the norm argument.

The simultaneous step below is stronger than constructing a separate
curve for each word.  Pairwise resultants have quadratic height cost
in the support, independent of the number of basis vectors.

\begin{theorem}[A common curve]
\label{p13:rig:common-curve}
There is an absolute constant $c>0$ with the following property.
Let $V$ be a space of even formal words of dimension $t\ge2$,
with a basis $w_1,\ldots,w_t$ satisfying
$|I_{w_i}|,|J_{w_i}|\le2q_0$, where $1\le q_0\le B$.
Suppose all $F_{w_i}(x,y)$ are integer squares and
$x,y\in[N,2N)$.  If $N>(B+1)^{cq_0^2}$, there is a
primitive irreducible polynomial $C_0\in\Z[X,Y]$, depending
on both variables, such that
\begin{equation}
 C_0(x,y)=0,\qquad \deg C_0\le48q_0,
 \qquad H(C_0)\le(B+1)^{cq_0}.
 \label{p13:rig:common-degree-height}
\end{equation}
In the function field $K_0=\mathbb Q(C_0=0)$, the image of
\[
               V\longrightarrow K_0^\times/K_0^{\times2},
                  \qquad w\longmapsto[F_w]
\]
has dimension at most one.  The constant is independent of $t$.
\end{theorem}

The proof is given in Section~\ref{p14:comp:rigidity:common-curve}.
Its key point is that every norm polynomial must share the same
irreducible factor through $(x,y)$: otherwise a resultant would give
a nonzero polynomial of height $(B+1)^{O(q_0^2)}$ with integer root
$x\ge N$.  Primitivity of the auxiliary coefficients then forces the
functional square-class image to have dimension at most one.

\subsection{Functional squares force integer fibres}

The centres have size $N$, but the next step counts points
in a box of side $O(B)$.  The bridge is that a square
function has even valuation at every point of the normalized
curve.  A simple zero belonging to a selected first-block
offset must therefore be matched by a selected second-block
offset.

\begin{proposition}[Integer fibres in the offset grid]
\label{p13:rig:grid-fibres}
Under the hypotheses of Theorem~\ref{p13:rig:common-curve},
assume both physical projections of $V$ are injective.
Put $D=\deg C_0$, $d_X=\deg_XC_0$, $d_Y=\deg_YC_0$.
For an absolute constant $c$,
\begin{equation}
 \#\{(u,v)\in[-B,B]^2\cap\Z^2:C_0(-u,-v)=0\}
       \ge\max(d_X,d_Y)(t-1-cD^2).
 \label{p13:rig:grid-count}
\end{equation}
A negative right side may be replaced by zero.
\end{proposition}

\begin{proof}
The kernel $V_0$ of the map to functional square classes
has dimension at least $t-1$.  The leading coefficients
and discriminants of $C_0$ in each variable are nonzero
polynomials, of total degree $O(D^2)$ in the remaining
variable.  Exclude the offsets at which one of these
polynomials vanishes.  Requiring the corresponding
coordinates of a word to be zero leaves a subspace $V_1$
of dimension at least $t-1-cD^2$.

Let $\mathcal U_X$ be the union of first-block offsets used
by $V_1$.  Projection injectivity gives
$|\mathcal U_X|\ge\dim V_1$.
For $u\in\mathcal U_X$, select a word containing $X+u$.
The noncritical fibre $X=-u$ has exactly $d_Y$ distinct
affine roots on the geometric normalization.  At each,
$X+u$ has a simple zero and no other first-block factor
vanishes.  The product for the selected word is a square
function, so a selected factor $Y+v$ must vanish there.
Every root is thus an integer $-v$ in the offset interval.
Distinct $u$ give disjoint fibres, proving at least
$d_Y\dim V_1$ grid points.  Interchanging $X,Y$ proves
the bound with $d_X$ as well.
The argument may be performed after extension to
$\mathbb C$ and normalization of the components; geometric
integrality is not required.
\end{proof}

\subsection{The point count at growing degree}

We use the integer-point part of Binyamini--Cluckers--Novikov
\cite[Theorem~2, equation~(27), in the specified author
version]{binyaminicluckersnovikov2024}.

\begin{theorem}[Binyamini--Cluckers--Novikov]
\label{p13:rig:bcn}
For an integral algebraic curve $\mathcal C\subset\mathbb A^n_K$
of degree $D$, where $n>1$ and $K\subset\mathbb C$, there are
a constant $c(n)$ and an absolute constant $\kappa$ such that
\[
 \#\{z\in\mathcal C\cap\Z^n:\max_i|z_i|\le H\}
       \le c(n)D^2H^{1/D}(\log H)^\kappa,
       \qquad H>2.
\]
The constants are independent of the curve's coefficients
and degree.  Integrality over $K$ is sufficient.
\end{theorem}

The proof in the Lean development gives the admissible value
$\kappa=17$. This logarithmic power need not be optimized here.

\begin{corollary}[Sublinear counts in the offset box]
\label{p13:rig:sublinear-grid}
Fix $0\le a<1/2$ and $c_0>0$. Let $D_{\max}=D_{\max}(B)$
be any fixed degree envelope such that, for every $\delta>0$,
$D_{\max}(B)\le B^{a+\delta}$ for all sufficiently large $B$.
Uniformly for integral plane curves of degree $2\le D\le D_{\max}(B)$,
\[
 \#(\mathcal C\cap[-c_0B,c_0B]^2\cap\Z^2)
 \ll_{c_0}(B^{1/2}+D_{\max}^2)(\log B)^\kappa=o(B).
\]
\end{corollary}

\begin{proof}
Set $H=\max(3,c_0B)$.  The logarithmic derivative of
$u^2H^{1/u}$ is $(2u-\log H)/u^2$, so its maximum on
$[2,M]$, $M\ge2$, is attained at an endpoint.
Put $M'=\max(M,\log H)$. Applying this observation on $[2,M']$,
and using $H^{1/M'}\le e$, gives, for every $2\le D\le M$,
\[
 D^2H^{1/D}
 \le 4H^{1/2}+e(M')^2
 \le4H^{1/2}+e\{M^2+(\log H)^2\}.
\]
The last logarithmic square is absorbed into $B^{1/2}$.
Apply Theorem~\ref{p13:rig:bcn} with $n=2$ and $M=D_{\max}(B)$.
Choose a fixed $\delta>0$ with $2(a+\delta)<1$. Both
$B^{1/2}(\log B)^\kappa$ and
$B^{2(a+\delta)}(\log B)^\kappa$ are $o(B)$.
The threshold is chosen after $c_0,a$, the envelope and the requested
error, and before the curve and its coefficients. No lower growth of
$D_{\max}$ is needed; bounded and polylogarithmic envelopes are included.
\end{proof}

This application uses the theorem's degree-uniform constant.
A bound with an unspecified constant depending on $D$ would
not give the conclusion of Corollary~\ref{p13:rig:sublinear-grid}.

\subsection{Nonaligned pairs have rank less than half a window}

\begin{theorem}[A uniform nonaligned rank bound]
\label{p13:rig:nonaligned}
For every fixed critical band,
\begin{equation}
 \limsup_{N\to\infty}\ 
 \sup_{\substack{x,y\in I_N\\x,y\ {\rm nonaligned}}}
             \frac{d_{xy}}B
       \le\frac23\log2<\frac12.
 \label{p13:rig:nonaligned-bound}
\end{equation}
Compatibility of the two events is not needed.
\end{theorem}

\begin{proof}
Suppose along a sequence that
$d_{xy}\ge((2/3)\log2+\delta)B$ for some fixed
$\delta>0$.  Choose a fixed $a<1/2$, sufficiently close to
$1/2$, that
$(2/3)\log(1/a)<(2/3)\log2+\delta/2$.
Theorem~\ref{p13:rig:short-basis} gives a space $V$
of dimension at least $\delta B/2-o(B)$ with injective
physical projections and a basis whose half-cardinalities
are at most $q_0=B^{a+o(1)}$.  The quantitative bridge to
the common-curve theorem is
\begin{equation}
              q_0^2\log(B+1)=B^{2a+o(1)}=o(B)=o(\log N).
 \label{p13:rig:threshold-bridge}
\end{equation}
Thus Theorem~\ref{p13:rig:common-curve} supplies an
irreducible $C_0$ through $(x,y)$ of degree
$D\le B^{a+o(1)}$ and logarithmic height $o(B)$, with
functional image of dimension at most one.

If $D\ge2$, Proposition~\ref{p13:rig:grid-fibres} gives
at least $\delta B/2-o(B)$ integer points in the offset
box, because $D^2=o(B)$.  Corollary~\ref{p13:rig:sublinear-grid}
gives $o(B)$ points, a contradiction.

If $D=1$, write $C_0=\alpha X+\beta Y+\gamma$.
Both $\alpha,\beta$ are nonzero.  Since
$H(C_0)=\exp(o(B))=N^{o(1)}$ and $x,y\ge N$, they have
opposite signs; otherwise $\gamma$ could not cancel the
two terms.  The line therefore has the form
$X=\lambda Y+\nu$ with $\lambda>0$.
The kernel of the functional-class map has positive
dimension.  Choose a nonzero word in it.  Each selection
is nonempty and even by projection injectivity.  On the
line, the distinct simple roots from the two blocks must
pair, each pairing giving $\nu+u=\lambda v$.
There are at least two pairings, hence
\[
 \lambda=\frac{u_1-u_2}{v_1-v_2}=\frac pq,
 \qquad (p,q)=1,\qquad 1\le p,q\le2B.
\]
Evaluating one pairing at the centres gives
$qx-py=pv-qu$, of absolute value at most $(p+q)B$.
This is an alignment as in \eqref{p13:rig:alignment},
contrary to the assumption.  Both degree cases are
impossible, and $\delta$ was arbitrary.
\end{proof}

The number $(2/3)\log2$ is a sufficient upper threshold
arising from the extraction and degree bounds.  No matching
rank construction is asserted for this constant.


\subsection{Aligned pairs and the exact contact distance}

For an aligned pair the line is already known.  Substitution into it
costs a height linear in the support, so the extraction exponent may
approach one without the quadratic resultant threshold.

\begin{proposition}[Rank on a known alignment]
\label{p13:rig:aligned}
Suppose $bx-ay=h$, where $(a,b)=1$, $1\le a,b\le B^3$,
and $|h|\le(a+b)B$.  For every fixed $\tau\in(0,1)$,
\begin{equation}
 d_{xy}\le\frac B{\max(a,b)}+\frac23B\log(1/\tau)
                    +O_{\tau,C}(B/\sqrt{\log B}).
 \label{p13:rig:aligned-quantitative}
\end{equation}
For distinct compatible centres this implies, uniformly
over aligned pairs,
$d_{xy}\le B/2+o_C(B)$.  If such a pair has an alignment
with $a=b=1$, then $d_{xy}=o_C(B)$.
\end{proposition}

The proof in Section~\ref{p14:comp:rigidity:aligned} counts the matched
linear factors on that line.  At height one, compatibility imposes
$|x-y|\ge L$; distance exactly $L=B-1$ leaves only one shared true
position.  A nonzero even relation needs at least two matches, so this
contact contributes no macroscopic intersection rank.  Thus the bound
includes the minimum compatible separation exactly.

\subsection{Uniform triple nullity and its probabilistic cost}

\begin{corollary}[Triple rank and maximum joint probability]
\label{p13:rig:triple}
Every compatible triple $T$ of distinct centres in the
fixed critical band satisfies
\begin{equation}
                 \rho_{\rm rel}(T)\le\frac32L+o_C(L).
 \label{p13:rig:triple-nullity}
\end{equation}
If any of its pairs is nonaligned, the stronger bound
$\rho_{\rm rel}(T)\le(1+(2/3)\log2)L+o_C(L)$ holds.
In particular, with $q(T)=\E\prod_{x\in T}J_x$,
\begin{equation}
                  \max_{\substack{T\subset I_N\\|T|=3}}q(T)
                            \le N^{-3/2+o_C(1)}.
 \label{p13:rig:max-triple}
\end{equation}
\end{corollary}

\begin{proof}
For three spaces of dimensions at most $L$, the dimension
identity and the intersection bound give
\[
 \begin{split}
 \rho_{\rm rel}(T)
 &=\dim(S_1\cap S_2)+\dim((S_1+S_2)\cap S_3)\\
 &\le L+\dim(S_1\cap S_2).
 \end{split}
\]
Permuting the indices gives
$\rho_{\rm rel}(T)\le L+\min_{i<j}d_{ij}$.
Every pair in a compatible triple is compatible.
Theorem~\ref{p13:rig:nonaligned} and
Proposition~\ref{p13:rig:aligned}, with $B=L+1$, prove
the stated nullity bounds.

For clarity, the passage to probabilities uses exact
affine constraints.  The prime signs induce a uniform
character on each finite square-class space.  The
constraints for $J_x=1$ have classes
$\{x-1,x\}$ and $\{x,x+j\}$, $1\le j<L$, with prescribed
values $-1$ and $+1$ respectively.  These selections span
$E_x$.  Thus an admissible window has
$p_x:=\E J_x=2^{-\dim S_x}$, and a compatible family has
\begin{equation}
 q(T)=2^{-\dim\sum_{x\in T}S_x}
                  =2^{\rho_{\rm rel}(T)}\prod_{x\in T}p_x.
 \label{p13:rig:exact-probability}
\end{equation}
Every local relation is supported on defective positions,
so Lemma~\ref{p13:rig:defective} implies
$L-O_C(B/\log B)\le\dim S_x\le L$.
Consequently $p_x=2^{-L+o_C(L)}$ for admissible windows;
inadmissible windows contribute zero.  Substitution in
\eqref{p13:rig:exact-probability} gives
$q(T)\le2^{-3L/2+o_C(L)}=N^{-3/2+o_C(1)}$ uniformly,
which is \eqref{p13:rig:max-triple}.
\end{proof}

The estimate \eqref{p13:rig:max-triple} is the arithmetic
input needed by the reconstruction from three-window
activities.  Its proof uses the true windows and exact
marginals, so no removal of locally exceptional windows
is required at this stage.
